% GENERATED FILE. Edit canonical lessons, then run npm run generate:books.
% canonical-source-hash: fnv1a64:4018bbff33a0de89
% canonical-lessons: MW-C90-dopar, MW-C90-athan, MW-C90-dinugai, MW-C90-parbhat, MW-C90-parsu

\chapter{Midday, Evening, In the early morning, Dawn, Day after tomorrow}
\label{ch:mw-midday-evening-in-the-early-morning-dawn-day-after-tomorrow}
\hlchaptermodality{90}

\begin{chapteropening}
\textbf{By the end of this chapter:} \emph{I can say midday, the evening, in the early morning, dawn and the day after tomorrow.}

Complete the last lesson of chapter 90: I can say midday, the evening, in the early morning, dawn and the day after tomorrow.
\end{chapteropening}

\section[\texorpdfstring{\mw{दोपार} (dopār)}{दोपार (dopār)}]{\mw{दोपार} (dopār) — midday, the afternoon}

\label{lesson:MW-C90-dopar}



\begin{quote}
Before the new one: say the Marwadi for a matter; talk, then the Marwadi for a language.

\emph{Read it:} \textbf{\mw{ये} \mw{कपड़ा} \mw{दिखावो।} \mw{ये} \mw{कितणे} \mw{कू} \mw{है}?} without stopping
\end{quote}

\subsection*{You'll want to know: \mw{दोपार}}
\textbf{\mw{दोपार}} — \emph{dopār} — \textquotedblleft{}midday, the afternoon\textquotedblright{}.

\begin{cousinweb}
\textbf{\mw{दोपार}} (\emph{dopār}) is usually explained as \emph{do}, \textquotedblleft{}two\textquotedblright{}, and \emph{pahar}, \textquotedblleft{}a watch of about three hours\textquotedblright{}, which CDIAL lists as Old Marwari \emph{pahara} under Sanskrit \emph{prahara}, \textquotedblleft{}a watch\textquotedblright{} (CDIAL 8900). CDIAL asks whether the Sanskrit word is itself a reshaping of Persian \emph{pahra}, \textquotedblleft{}watch\textquotedblright{}. It has no entry for \textbf{\mw{दोपार}} (\emph{dopār}).
\end{cousinweb}

\begin{grammarlens}[title={What you've built}]
One more word for this chapter.
\end{grammarlens}

\subsection*{Your turn}
\begin{itemize}
\raggedright
  \item \emph{Say it:} \emph{dopār}
  \item \emph{Say it:} \emph{dopār}, once more
  \item \emph{Say it:} say \emph{dopār}
  \item \emph{Recall:} say the Marwadi for a matter; talk, then the Marwadi for a language, then say \emph{dopār} again
\end{itemize}

\begin{tcolorbox}[breakable,colback=teal!4,colframe=teal!35!black,title={Before you move on}]
What does \mw{दोपार} mean? (\textquotedblleft{}Midday, the afternoon\textquotedblright{}.) Say it once more.
\end{tcolorbox}
\section[\texorpdfstring{\mw{आथण} (āthaṇ)}{आथण (āthaṇ)}]{\mw{आथण} (āthaṇ) — the evening}

\label{lesson:MW-C90-athan}



\begin{quote}
Before the new one: say the Marwadi for a language, then the Marwadi for midday.
\end{quote}

\subsection*{You'll want to know: \mw{आथण}}
\textbf{\mw{आथण}} — \emph{āthaṇ} — \textquotedblleft{}the evening\textquotedblright{}.

\begin{cousinweb}
\textbf{\mw{आथण}} (\emph{āthaṇ}) is a tadbhava, a word inherited from Sanskrit and reshaped by sound change. It comes from \emph{astamayana}, \textquotedblleft{}setting of the sun\textquotedblright{}. CDIAL lists Old Marwari \emph{āthaṇa} under it (CDIAL 975). It is built on \emph{astam eti}, \textquotedblleft{}goes home\textquotedblright{}, said of the sun.
\end{cousinweb}

\begin{grammarlens}[title={What you've built}]
One more word for this chapter.
\end{grammarlens}

\subsection*{Your turn}
\begin{itemize}
\raggedright
  \item \emph{Say it:} \emph{āthaṇ}
  \item \emph{Say it:} \emph{āthaṇ}, once more
  \item \emph{Say it:} say \emph{āthaṇ}
  \item \emph{Recall:} say the Marwadi for a language, then the Marwadi for midday, then say \emph{āthaṇ} again
\end{itemize}

\begin{tcolorbox}[breakable,colback=teal!4,colframe=teal!35!black,title={Before you move on}]
What does \mw{आथण} mean? (\textquotedblleft{}The evening\textquotedblright{}.) Say it once more.
\end{tcolorbox}
\section[\texorpdfstring{\mw{दिनूगै} (dinūgai)}{दिनूगै (dinūgai)}]{\mw{दिनूगै} (dinūgai) — in the early morning}

\label{lesson:MW-C90-dinugai}



\begin{quote}
Before the new one: say the Marwadi for midday, then the Marwadi for the evening.
\end{quote}

\subsection*{You'll want to know: \mw{दिनूगै}}
\textbf{\mw{दिनूगै}} — \emph{dinūgai} — \textquotedblleft{}in the early morning\textquotedblright{}.

\begin{cousinweb}
CDIAL has no entry for \textbf{\mw{दिनूगै}} (\emph{dinūgai}), so its history is not settled. It seems to be \emph{din}, \textquotedblleft{}day\textquotedblright{} (CDIAL 6328), with a form of the verb CDIAL lists as Old Marwari \emph{ūgaï}, \textquotedblleft{}rises, appears\textquotedblright{}, under \emph{*udgāti}, \textquotedblleft{}rises\textquotedblright{} (CDIAL 1954): the time when the day comes up.
\end{cousinweb}

\begin{grammarlens}[title={What you've built}]
One more word for this chapter.
\end{grammarlens}

\subsection*{Your turn}
\begin{itemize}
\raggedright
  \item \emph{Say it:} \emph{dinūgai}
  \item \emph{Say it:} \emph{dinūgai}, once more
  \item \emph{Say it:} say \emph{dinūgai}
  \item \emph{Recall:} say the Marwadi for midday, then the Marwadi for the evening, then say \emph{dinūgai} again
\end{itemize}

\begin{tcolorbox}[breakable,colback=teal!4,colframe=teal!35!black,title={Before you move on}]
What does \mw{दिनूगै} mean? (\textquotedblleft{}In the early morning\textquotedblright{}.) Say it once more.
\end{tcolorbox}
\section[\texorpdfstring{\mw{परभात} (parbhāt)}{परभात (parbhāt)}]{\mw{परभात} (parbhāt) — dawn}

\label{lesson:MW-C90-parbhat}



\begin{quote}
Before the new one: say the Marwadi for the evening, then the Marwadi for in the early morning.
\end{quote}

\subsection*{You'll want to know: \mw{परभात}}
\textbf{\mw{परभात}} — \emph{parbhāt} — \textquotedblleft{}dawn\textquotedblright{}.

\begin{cousinweb}
\textbf{\mw{परभात}} (\emph{parbhāt}) is a semi-tatsama: Sanskrit \emph{prabhāta}, \textquotedblleft{}shone forth; daybreak\textquotedblright{}, borrowed after the old sound changes had run their course and then lightly reshaped (CDIAL 8707). Marwadi splits the \emph{pr} with a vowel. The inherited line shows in Prakrit \emph{pahāya}, \textquotedblleft{}dawn\textquotedblright{}, which CDIAL lists under it.
\end{cousinweb}

\begin{grammarlens}[title={What you've built}]
One more word for this chapter.
\end{grammarlens}

\subsection*{Your turn}
\begin{itemize}
\raggedright
  \item \emph{Say it:} \emph{parbhāt}
  \item \emph{Say it:} \emph{parbhāt}, once more
  \item \emph{Say it:} say \emph{parbhāt}
  \item \emph{Recall:} say the Marwadi for the evening, then the Marwadi for in the early morning, then say \emph{parbhāt} again
\end{itemize}

\begin{tcolorbox}[breakable,colback=teal!4,colframe=teal!35!black,title={Before you move on}]
What does \mw{परभात} mean? (\textquotedblleft{}Dawn\textquotedblright{}.) Say it once more.
\end{tcolorbox}
\section[\texorpdfstring{\mw{परसूं} (parsū̃)}{परसूं (parsū̃)}]{\mw{परसूं} (parsū̃) — the day after tomorrow; the day before yesterday}

\label{lesson:MW-C90-parsu}



\begin{quote}
Before the new one: say the Marwadi for in the early morning, then the Marwadi for dawn.
\end{quote}

\subsection*{You'll want to know: \mw{परसूं}}
\textbf{\mw{परसूं}} — \emph{parsū̃} — \textquotedblleft{}the day after tomorrow; the day before yesterday\textquotedblright{}.

\begin{cousinweb}
\textbf{\mw{परसूं}} (\emph{parsū̃}) is a tadbhava, a word inherited from Sanskrit and reshaped by sound change. It comes from \emph{paraśvas}, \textquotedblleft{}the day after tomorrow\textquotedblright{}. CDIAL gives no Marwari form in this entry (CDIAL 7799j); the match is through Hindi \emph{parsū̃}. It is built on \emph{para}, \textquotedblleft{}beyond\textquotedblright{}, and \emph{śvas}, \textquotedblleft{}tomorrow\textquotedblright{}; like Hindi, Marwadi uses it for both directions in time.
\end{cousinweb}

\begin{grammarlens}[title={What you've built}]
One more word for this chapter.
\end{grammarlens}

\subsection*{Your turn}
\begin{itemize}
\raggedright
  \item \emph{Say it:} \emph{parsū̃}
  \item \emph{Say it:} \emph{parsū̃}, once more
  \item \emph{Say it:} say \emph{parsū̃}
  \item \emph{Recall:} say the Marwadi for in the early morning, then the Marwadi for dawn, then say \emph{parsū̃} again
\end{itemize}

\begin{tcolorbox}[breakable,colback=teal!4,colframe=teal!35!black,title={Before you move on}]
What does \mw{परसूं} mean? (\textquotedblleft{}The day after tomorrow; the day before yesterday\textquotedblright{}.) Say it once more.
\end{tcolorbox}
